% Overview: `linear-pipeline` (design-axes.md axis 15; the generic form, and
% what P1's training-loop figure reduces to without its artwork).
% Four stages left to right, each led by its role's icon; a tint over the
% stages this paper contributed and a grey one over what every compared arm
% already had; frozen / trained as icon badges rather than words; a legend that
% keys both. Carries no magnitudes.
%
% Icons come from the five role macros (\acIconInput and friends, preamble.tex).
% A design draws them; if the drawn glyph is wrong for this paper's domain,
% renew it once in macros.tex from references/figure-icons.md rather than
% hard-coding a name here.
%
% ------------------------------------------------------------ density pass
% The four-stage row above is the whole figure only when the mechanism has
% nothing repeated for a reader to see happen. Most methods do: a layer, a
% timestep, a sample, an iteration, a candidate -- some countable sub-unit
% the mechanism treats differently from a baseline. When your paper has one,
% open it up below the row as a WORKED INSTANCE, the way this grammar's own
% proof does (paper-batch-v2/S10-fewshot-segmentation, PromptSeg: which of
% the encoder's 12 blocks receive a prompt). That pass took this grammar
% from 18 drawn objects to 51 on a page-width float, and the structure that
% kept it legible rather than noisy was:
%
%   - ONE ruler of positions/items, reused by every row below it. Never
%     redraw the ruler per row -- that is what makes two rows a comparison
%     instead of two unrelated diagrams.
%   - a cell's STATUS is its own fill and glyph (`acitem*` below), never a
%     neutral cell plus a badge beside it.
%   - if your paper's rule is a decision (a threshold, a criterion, a
%     learned gate) rather than a fixed placement, add the small
%     feature-to-decision cluster at the bottom: `acfeat` nodes for 2-5
%     things the decision reads, feeding one small node (a spare `accard`
%     works) naming the rule. Delete it if the mechanism makes no per-item
%     decision.
%   - a bracket under a row's active span carries a SYMBOL, never a count
%     read off a result (references/figure-ladder.md, Tier C): an
%     architecture constant such as a fixed $N$ is fine to write in full,
%     the way \Cref{eq:inject}'s own $l=1,\dots,N$ is, because it is not a
%     measured effect.
%
% If your paper has no such repeated sub-unit -- a one-shot transform with
% nothing to index -- leave the worked instance out rather than inventing
% one; the four-stage row alone is a complete, honest figure, just a
% shorter one.
%
% \acitem / acitemkeep / acitemdrop / acitemgate / acitemna, \acruler,
% \acrowlabel and \acfeat are declared once in preamble.tex's shared
% tikzset block (not here), so any linear-pipeline figure reuses them
% without redeclaring anything. Anchor every new row to the ALREADY-PLACED
% ruler node above it (`below=... of l1`, never a hand-computed coordinate
% guessed before the row above it was drawn) -- a coordinate that exists but
% whose true extent does not yet is exactly the bug that made an earlier
% draft's callout print underneath its own section header.
%
% source-data: none (Tier C structural diagram; see references/figure-ladder.md)
\begin{acwidefigure}
  \centering
  \acfit{%
  \begin{tikzpicture}[x=1cm,y=1cm]
    \node[accard,text width=2.5cm] (a) at (0,0)
      {\acicon{\acIconInput}\\[1pt]\textbf{\slot{Input}}\\\slot{what enters}};
    \node[actrap,right=0.75 of a,text width=2.2cm] (b)
      {\acicon{\acIconModel}\\[1pt]\textbf{\slot{Shared stage}}\\\slot{common to every arm}};
    \node[accard,right=1.05 of b,text width=2.5cm] (c)
      {\acicon{\acIconMech}\\[1pt]\textbf{\slot{The mechanism}}\\\slot{what this paper adds}};
    \node[accard,right=0.75 of c,text width=2.5cm] (d)
      {\acicon{\acIconOut}\\[1pt]\textbf{\slot{Output}}\\\slot{what is reported}};
    % the representation between the shared stage and the mechanism, drawn
    % rather than named: six tokens, no magnitude implied.
    \node[inner sep=0pt] (rep) at ($(b.east)!0.5!(c.west)$) {\actokens{6}};
    \node[actag,text=ACIcon,above=1pt of rep] {\slot{$h$}};
    \draw[acarrow] (a) -- (b);
    \draw[acarrow] (b) -- (rep);
    \draw[acarrow] (rep) -- (c);
    \draw[acarrow] (c) -- (d);
    % State badges go BELOW the cards. Above, they collided with the region
    % tag, which sits on the tinted box's top edge: two small uppercase
    % labels at the same height read as one broken line.
    \node[actag,below=3pt of b] (bf) {\acsnow\ \slot{FROZEN}};
    \node[actag,below=3pt of c,text=ACBlue] (ct) {\acflame\ \slot{TRAINED}};
    \begin{scope}[on background layer]
      \node[acbase,fit=(a)(b)(bf),inner sep=8pt] (base) {};
      \node[acours,fit=(c)(d)(ct),inner sep=8pt] (ours) {};
    \end{scope}
    \node[actag,anchor=south west] at ([xshift=3pt,yshift=2pt]base.north west)
      {\slot{COMMON TO EVERY ARM}};
    \node[actag,anchor=south west,text=ACBlue] at ([xshift=3pt,yshift=2pt]ours.north west)
      {\slot{CONTRIBUTED BY THIS PAPER}};
    \node[aclegend,anchor=north west] (toplegend) at ($(base.south west)+(0,-0.35)$)
      {\acsnow~\slot{what is not updated} \quad \acflame~\slot{what is}\\[2pt]
       \textcolor{ACBlue!80}{\rule[0.12em]{5mm}{0.8pt}}~\slot{what flows along the solid arrows}};
    \node[acsoft,text width=4.6cm,align=center,anchor=north west]
      (callout) at ($(ours.south west)+(0.2,-0.35)$)
      {\slot{callout: the one design decision that makes the comparison fair}};

    % ================================================================
    % WORKED INSTANCE (delete this whole block if the mechanism has no
    % repeated sub-unit to index -- see the note above). \slot{N} is
    % whatever your method counts: layers, steps, samples, candidates.
    % The ruler below shows a representative window with an ellipsis; if
    % your $N$ is small enough to show in full (5 or fewer), delete the
    % ellipsis node and draw every position instead.
    % ================================================================
    \coordinate (bottommid) at ($(toplegend.south)!0.5!(callout.south)$);
    \node[actag,text width=8.6cm,align=center,anchor=north,
          below=0.5 of bottommid] (zoomhdr)
      {\slot{WHAT HAPPENS AT EACH OF THE $N$ \dots (SAMPLES / STEPS / LAYERS)} ---
       \slot{SAME INPUT, TWO RULES BELOW}};

    \node[acruler,below=0.7 of zoomhdr,xshift=-3.4cm] (i1) {\slot{$1$}};
    \node[acruler,right=0.55 of i1]                   (i2) {\slot{$2$}};
    \node[acruler,right=0.55 of i2]                   (idots) {$\cdots$};
    \node[acruler,right=0.55 of idots]                (in1) {\slot{$N{-}1$}};
    \node[acruler,right=0.55 of in1]                  (in)  {\slot{$N$}};

    % Row A: \slot{the baseline or restricted condition}. Swap acitemgate /
    % acitemna for acitemkeep / acitemdrop if the row's status is a
    % correctness judgment (kept/dropped) rather than a presence/absence
    % one -- whichever pair matches what your mechanism actually decides.
    \node[acitemgate,below=0.5 of i1]     (rA1) {\slot{status glyph or blank}};
    \node[acitemgate,below=0.5 of i2]     (rA2) {\slot{status glyph or blank}};
    \node[acitemna,  below=0.5 of idots]  (rAdots) {\slot{status glyph or blank}};
    \node[acitemna,  below=0.5 of in1]    (rAn1) {\slot{status glyph or blank}};
    \node[acitemna,  below=0.5 of in]     (rAn) {\slot{status glyph or blank}};
    \node[acrowlabel,left=0.35 of rA1] (rowA) {\slot{baseline rule}};
    \draw[ACBlue!65,line width=0.7pt]
      (rA1.south west)+(0,-0.08) -- ++(0,-0.08)
      -- ($(rA2.south east)+(0,-0.16)$) -- ++(0,0.08);
    \node[actag,text=ACBlue,anchor=north] at ($(rA1.south east)!0.5!(rA2.south west)+(0,-0.34)$)
      {\small\slot{span, symbol only}};

    % Row B: \slot{this paper's own rule}. Every cell is `below=1.0 of` its
    % row-A counterpart -- a full centimetre of real clearance under row A's
    % own bracket and label, not a hand-picked offset.
    \node[acitemgate,below=1.0 of rA1]    (rB1) {\slot{status glyph or blank}};
    \node[acitemgate,below=1.0 of rA2]    (rB2) {\slot{status glyph or blank}};
    \node[acitemgate,below=1.0 of rAdots] (rBdots) {\slot{status glyph or blank}};
    \node[acitemgate,below=1.0 of rAn1]   (rBn1) {\slot{status glyph or blank}};
    \node[acitemgate,below=1.0 of rAn]    (rBn) {\slot{status glyph or blank}};
    \node[acrowlabel,left=0.35 of rB1] (rowB) {\slot{this paper's rule}};
    \draw[ACBlue!65,line width=0.7pt]
      (rB1.south west)+(0,-0.08) -- ++(0,-0.08)
      -- ($(rBn.south east)+(0,-0.16)$) -- ++(0,0.08);
    \node[actag,text=ACBlue,anchor=north] at ($(rB1.south east)!0.5!(rBn.south west)+(0,-0.34)$)
      {\small\slot{span, symbol only}};

    % ----------------------------------------------------------------
    % OPTIONAL: what the mechanism reads to make its per-item call. Delete
    % this cluster (and its two connecting draws below) if the rule above
    % is a fixed placement rather than a learned or thresholded decision.
    % ----------------------------------------------------------------
    \node[acfeat,below=1.05 of rB1,xshift=1.3cm] (f1) {\slot{feature 1}};
    \node[acfeat,right=0.5 of f1]                (f2) {\slot{feature 2}};
    \node[actag,above=3pt of f1,anchor=south west,text=ACIcon,text width=3.4cm]
      {\slot{what these features are, cite the equation}};
    \draw[acstat] (rB1.south) .. controls +(0,-0.7) and +(-0.3,0.5) .. (f1.north);
    \node[accard,text width=3.1cm,font=\tiny,below=0.7 of f1,xshift=0.6cm] (decide)
      {\aciconat{ACIcon}{\acIconMech}\ \slot{the rule this paper's mechanism applies to f1, f2, ...}};
    \draw[acstat] (f1.south) -- (decide.north west);
    \draw[acstat] (f2.south) -- (decide.north east);

    % ----------------------------------------------------------------
    % Legend: keys every encoding introduced below the top row. The
    % frozen/trained line is already keyed in `toplegend` above; do not
    % repeat it here unless this legend has drifted out of that one's reach.
    % ----------------------------------------------------------------
    \node[aclegend,below=0.55 of decide,text width=10.4cm,anchor=north]
      {\slot{status glyph}~\slot{what it means at a position} \quad
       \slot{status glyph}~\slot{what the other state means}\\[3pt]
       \textcolor{ACBlue!70}{\rule[0.15em]{2.5mm}{0.8pt}\,\rule[0.15em]{2.5mm}{0.8pt}}~\slot{what the span bracket marks, under each rule}};
  \end{tikzpicture}}
  \caption{\textbf{\slot{name the pipeline}.} \slot{Top}: \slot{what is shown, what is
    compared, and what to notice, in two or three sentences}. \slot{Bottom}:
    \slot{one or two sentences on what the worked instance shows at each
    position, and what to notice comparing the two rows}.}
  \label{fig:overview}
\end{acwidefigure}
